\section{Conclusão}

O plano é coerente porque cresce na ordem em que reduz incertezas: esclarecer a experiência, testar uso autônomo, garantir oferta local, medir retorno e só então validar cobrança e ampliar território. A literatura sobre decisão empreendedora baseada em hipóteses [4], plataformas de dois lados [9, 10] e escala prematura [8] reforça essa sequência. O setor pet brasileiro oferece contexto relevante [19], mas não comprova demanda capturável. O resultado esperado desta fase é evidência operacional confiável: tutores ativados e retidos, parceiros que retornam, custos conhecidos e limites clínicos e de dados respeitados.
